% Chapter 13: Mind-Change Ordinals and Transfinite Hierarchies
%   mind_change_ordinal = Profile C (revelation): narrate the Freivalds-Smith discovery
%   constructive_ordinal = Profile A (vocabulary): brief setup of notation
%   bc_learning = Profile D (negative result): BC \ EX separation witness is the content
%   ordinal hierarchy = Profile B (load-bearing): full proof of strict separation
%   mind_change_bound = operational, brief
%   anomalous_learning = definition + hierarchy, brief
%
% Texture: recursion-theoretic with ordinal arithmetic. This chapter extends
% the diagonalization methods of Ch7 into the transfinite. The surprise is
% that ordinals are not merely convenient bookkeeping but the *right* measure.

\chapter{Mind-Change Ordinals and Transfinite Hierarchies}
\label{ch:ordinals}

How many times must a learner change its mind before converging to a correct
hypothesis?  In \Cref{ch:gold} we saw that $\Ex$-learners are permitted
finitely many mind changes, but the number may be unbounded: no single
integer $n$ suffices to bound the mind changes across all target functions in
the class.  This observation raises a natural question: can one \emph{measure}
the mind-change complexity of a class, and does that measure yield a strict
hierarchy of learning power?

The answer, due to Freivalds and Smith~\cite{FreivaldsSmith1993}, is that
the correct measure is not an integer but a \emph{constructive ordinal}.
The mind-change ordinal of a class assigns to each learnable class a
position in a transfinite hierarchy that refines the gap between finite
learning ($\FIN$) and identification in the limit ($\Ex$).  The hierarchy
is strict at every level, and its union does not exhaust $\Ex$.

This chapter develops the ordinal mind-change hierarchy in full.  We begin
with the minimal ordinal vocabulary needed (\Cref{sec:constructive-ordinals}),
then tell the story of how ordinals entered learning theory
(\Cref{sec:mco-revelation}), prove that the hierarchy is strict
(\Cref{sec:ordinal-hierarchy}), and close with the separation between
behaviorally correct and explanatory learning (\Cref{sec:bc-separation}).

\medskip

\begin{enumerate}[leftmargin=*,itemsep=2pt]
\item \textbf{Constructive ordinals} (\Cref{sec:constructive-ordinals}):
  Kleene's system~$\mathcal{O}$, notation for computable ordinals.
\item \textbf{The mind-change ordinal} (\Cref{sec:mco-revelation}):
  how ordinals became the right measure of learning complexity.
\item \textbf{The strict hierarchy} (\Cref{sec:ordinal-hierarchy}):
  $\forall \alpha < \beta$, $\Ex_\alpha \subsetneq \Ex_\beta$, with full proof.
\item \textbf{Anomalous learning} (\Cref{sec:anomalous}): the $\Ex_a$
  hierarchy and its interaction with mind changes.
\item \textbf{Behaviorally correct learning} (\Cref{sec:bc-separation}):
  $\BC \setminus \Ex$ via diagonal witness.
\item \textbf{The full landscape} (\Cref{sec:ordinal-landscape}):
  a diagram and summary of all inclusions.
\end{enumerate}


% ================================================================
% SECTION 1: CONSTRUCTIVE ORDINALS (Profile A -- vocabulary)
% ================================================================

\section{Constructive Ordinals}
\label{sec:constructive-ordinals}

We need ordinals that a Turing machine can manipulate.  The full class of
countable ordinals is too large: many have no computable description.  Kleene's
system~$\mathcal{O}$ provides the computable fragment.

\begin{defn}[Kleene's System $\mathcal{O}$]
\label{def:system-O}
The \emph{constructive ordinals} are a subset $\mathcal{O} \subseteq \N$
together with a partial ordering $<_{\mathcal{O}}$ and a map
$|\cdot| : \mathcal{O} \to \omega_1^{\mathrm{CK}}$ (the Church--Kleene
ordinal) defined inductively:
\begin{enumerate}[leftmargin=*,itemsep=2pt]
\item $1 \in \mathcal{O}$ with $|1| = 0$.
\item If $a \in \mathcal{O}$, then $2^a \in \mathcal{O}$ with
  $|2^a| = |a| + 1$, and $a <_{\mathcal{O}} 2^a$.
\item If $\varphi_e$ is a total computable function with
  $\varphi_e(0) <_{\mathcal{O}} \varphi_e(1) <_{\mathcal{O}} \cdots$,
  then $3 \cdot 5^e \in \mathcal{O}$ with
  $|3 \cdot 5^e| = \sup_n |\varphi_e(n)|$, and
  $\varphi_e(n) <_{\mathcal{O}} 3 \cdot 5^e$ for all $n$.
\end{enumerate}
\end{defn}

\begin{notation}
\label{not:ordinal-shorthand}
We write ordinals in standard notation ($0, 1, 2, \ldots, \omega,
\omega + 1, \ldots, \omega \cdot 2, \ldots, \omega^2, \ldots,
\omega^\omega, \ldots$) rather than their codes in $\mathcal{O}$.
When we say ``ordinal $\alpha$'' in the context of mind-change bounds,
we always mean a constructive ordinal: one that has a notation in
$\mathcal{O}$.
\end{notation}

\begin{remark}
\label{rem:system-O-noncomputable}
Membership in $\mathcal{O}$ is $\Pi_1^1$-complete and hence far from
decidable.  This does not obstruct its use as a mind-change counter:
the learner is given a specific notation $a \in \mathcal{O}$ and needs
only to compute the predecessor relation, which \emph{is} computable
given the notation.  The learner never needs to decide whether an
arbitrary number belongs to $\mathcal{O}$.
\end{remark}


% ================================================================
% SECTION 2: THE MIND-CHANGE ORDINAL (Profile C -- revelation)
% ================================================================

\section{The Mind-Change Ordinal: How Ordinals Entered Learning Theory}
\label{sec:mco-revelation}

The story begins with a natural frustration.

\subsection*{The problem that was stuck}

After Gold's 1967 theorem established the basic framework and Case and
Smith's 1983 paper~\cite{CaseSmith1983} introduced the hierarchy of
success criteria ($\FIN$, $\Ex$, $\BC$, and their variants), a question
remained open throughout the 1980s: \emph{how should one measure the
mind-change complexity of a class?}

The naive approach is to count mind changes with natural numbers.  Define
$\Ex_n$ as the collection of classes learnable by a machine that changes
its mind at most $n$ times on any input.  This gives a hierarchy:
\[
  \FIN = \Ex_0 \subsetneq \Ex_1 \subsetneq \Ex_2 \subsetneq \cdots
  \subsetneq \Ex.
\]
Velauthapillai~\cite{Velauthapillai1989} proved these inclusions are
strict.  But the hierarchy has a defect: the union
$\bigcup_{n \in \N} \Ex_n$ does \emph{not} equal $\Ex$.

\begin{defn}[Mind-Change Count]
\label{def:mc-count}
Let $M$ be a learner and $f$ a target function.  The \emph{mind-change
count} of $M$ on $f$, denoted $\mathrm{mc}(M, f)$, is the number of
times $M$ changes its hypothesis while processing a text for $f$:
\[
  \mathrm{mc}(M, f) \;=\; \bigl|\{t \in \N : h_t \neq h_{t+1}\}\bigr|.
\]
The \emph{mind-change complexity} of a class $\mathcal{C}$ with respect
to $M$ is $\mathrm{mc}(M, \mathcal{C}) = \sup_{f \in \mathcal{C}}
\mathrm{mc}(M, f)$.
\end{defn}

The problem is that $\mathrm{mc}(M, \mathcal{C})$ may be finite for every
$f \in \mathcal{C}$ yet unbounded across the class.  The supremum is
$\omega$, but no single integer bounds all instances.  There exist
$\Ex$-learnable classes in $\Ex \setminus \bigcup_n \Ex_n$: the learner
converges on every input but there is no uniform bound on the number of
mind changes.

\begin{computational}
Consider the class $\mathcal{C} = \{f_n : n \in \N\}$ where $f_n(x) = 0$
for $x < n$ and $f_n(x) = 1$ for $x \geq n$.  A learner that guesses
``$f_n$ for the largest $n$ seen so far'' is an $\Ex$-learner for
$\mathcal{C}$.  On input $f_n$, it changes its mind exactly $n$ times
(from $f_0$ to $f_1$ to $\ldots$ to $f_n$).  The mind-change count is
$n$, which is finite for each $f_n$ but unbounded across $\mathcal{C}$.
No $\Ex_k$ contains $\mathcal{C}$.
\end{computational}

\subsection*{The obvious approach that fails}

One might try to rescue the integer hierarchy by allowing $\Ex_\omega$
as a ``catch-all'' for classes learnable with finitely many but unbounded
mind changes.  But this merely relabels the gap; it does not give a
\emph{refined} hierarchy.  The gap between $\bigcup_n \Ex_n$ and $\Ex$
may itself contain structure: some classes might require the learner
to perform mind changes that depend on mind changes, a kind of
second-order revision.  Is there a hierarchy that captures this structure?

\subsection*{The Freivalds--Smith construction}

In 1993, Freivalds and Smith~\cite{FreivaldsSmith1993} answered this
question with a construction that remains surprising.  They observed
that the \emph{ordinal arithmetic of mind changes} is the correct
framework.

\begin{defn}[Ordinal Mind-Change Bound]
\label{def:ordinal-mc}
Let $\alpha$ be a constructive ordinal.  A learner $M$
\emph{$\alpha$-bounds its mind changes} if $M$ carries a counter
initialized to some notation for $\alpha$, and at each mind change,
$M$ replaces the counter with a strictly smaller ordinal
(in $<_{\mathcal{O}}$).  Since there is no infinite descending
sequence in the ordinals, $M$ must eventually stop changing its mind.
\end{defn}

\begin{defn}[$\Ex_\alpha$]
\label{def:ex-alpha}
For a constructive ordinal $\alpha$, the class $\Ex_\alpha$ consists
of all concept classes $\mathcal{C}$ such that some learner
$\alpha$-bounds its mind changes and $\Ex$-identifies $\mathcal{C}$.
\end{defn}

\begin{remark}[The hierarchy is notation-relative]
\label{rem:notation-dependence}
A subtlety underlies every statement that follows.  The class
$\Ex_\alpha$ is defined through a \emph{notation} $a \in \mathcal{O}$
for~$\alpha$, not through the ordinal $|a|$ alone, because the learner
decrements along the computable predecessor sequence that the notation
supplies.  Two notations $a, a'$ with $|a| = |a'|$ can give
$\Ex_a \neq \Ex_{a'}$ (Exercise~6).  When the theorems below write
``$\Ex_\alpha$'' and ``for all constructive~$\alpha$,'' they mean: fix a
notation for each ordinal; the strictness statements hold for these
fixed notations.  Which ordinals (as opposed to notations) admit a
notation-invariant $\Ex_\alpha$ is the subject
of~\cite{MartinSharmaStephan2006} and is only partially understood.
\end{remark}

The key insight is operational: a counter initialized to $\omega$ can
be decremented to any finite value $n$, but a counter initialized to
$n$ can only be decremented $n$ times.  So $\Ex_\omega$ is strictly
stronger than every $\Ex_n$: a learner with an $\omega$-counter can
``spend'' its first mind change moving from $\omega$ to some finite $n$,
after which it has $n$ mind changes remaining.  The value $n$ can
depend on the data seen so far.

\begin{historical}
Freivalds and Smith's construction was motivated by an analogy with
proof theory, where ordinal analysis assigns ordinal ``proof-theoretic
strengths'' to formal systems.  Just as the consistency strength of
Peano arithmetic is measured by the ordinal $\varepsilon_0$, the
mind-change complexity of a class is measured by a constructive ordinal.
The analogy is not perfect (the ordinals in learning theory stay well
below $\varepsilon_0$), but the structural parallel guided the discovery.
\end{historical}

\begin{example}
\label{ex:omega-learner}
Recall the class $\mathcal{C} = \{f_n : n \in \N\}$ from the
computational illustration above.  A learner with an $\omega$-counter
operates as follows:
\begin{enumerate}[leftmargin=*,itemsep=2pt]
\item Initialize the counter to $\omega$.
\item Upon seeing the first datum $f(0)$, make an initial guess and
  set the counter to some finite value (say, the current time step).
\item Each subsequent mind change decrements the (now finite) counter.
\end{enumerate}
After the first mind change, the learner has a finite counter, so it
changes its mind at most finitely many times.  The key point is
that the size of the finite counter is chosen \emph{adaptively}: it
depends on the data.  This is exactly what an $\omega$-counter provides
and a fixed integer counter does not.
\end{example}

\begin{remark}
\label{rem:ordinal-vs-integer}
The passage from $\bigcup_n \Ex_n$ to $\Ex_\omega$ is
not a mere notational trick.  It reflects a genuine increase in
computational power: the learner can defer its commitment to a finite
bound until it has seen enough data to determine what the bound should
be.  This deferred-commitment mechanism is the operational content of
transfinite mind changes.
\end{remark}


% ================================================================
% SECTION 3: THE STRICT HIERARCHY (Profile B -- load-bearing)
% ================================================================

\section{The Ordinal Mind-Change Hierarchy}
\label{sec:ordinal-hierarchy}

The central theorem of this chapter asserts that the ordinal-indexed
hierarchy is strict at every level.

\begin{thm}[Freivalds--Smith 1993; Ambainis--Jain--Sharma 1999]
\label{thm:ordinal-hierarchy}
For all constructive ordinals $\alpha < \beta$,
\[
  \Ex_\alpha \subsetneq \Ex_\beta \subsetneq \Ex.
\]
Moreover, $\bigcup_{\alpha \in \mathcal{O}} \Ex_\alpha \subsetneq \Ex$.
\end{thm}

The proof has two components: the inclusion $\Ex_\alpha \subseteq
\Ex_\beta$ (easy), and the strictness $\Ex_\alpha \neq \Ex_\beta$
(which requires constructing a witness class in $\Ex_\beta \setminus
\Ex_\alpha$).  We develop each in turn.

\subsection{Inclusions}

\begin{lemma}
\label{lem:inclusion}
If $\alpha <_{\mathcal{O}} \beta$, then $\Ex_\alpha \subseteq \Ex_\beta$.
\end{lemma}

\begin{proof}
Let $\mathcal{C} \in \Ex_\alpha$, witnessed by learner $M$ with an
$\alpha$-counter.  Since $\alpha < \beta$, we can initialize a
$\beta$-counter and immediately decrement it to $\alpha$.  (Formally:
replace the counter $\beta$ with $\alpha$ at the first step, which is
a valid decrement since $\alpha <_{\mathcal{O}} \beta$.)  Then run $M$
with the remaining $\alpha$-counter.  This witnesses
$\mathcal{C} \in \Ex_\beta$.
\end{proof}

\subsection{Witness classes for strictness}

The hard direction requires, for each pair $\alpha < \beta$, a class
$\mathcal{C}_{\alpha,\beta} \in \Ex_\beta \setminus \Ex_\alpha$.  We
construct this via a \emph{self-referential coding} argument.

\begin{lemma}[Witness Construction]
\label{lem:witness}
For every constructive ordinal $\alpha$, there exists a class
$\mathcal{C}_\alpha \in \Ex_{\alpha+1} \setminus \Ex_\alpha$.
\end{lemma}

\begin{proof}
We construct $\mathcal{C}_\alpha$ so that:
\begin{enumerate}[label=(\roman*)]
\item an $(\alpha+1)$-bounded learner can identify every function
  in $\mathcal{C}_\alpha$, but
\item no $\alpha$-bounded learner can identify all functions in
  $\mathcal{C}_\alpha$.
\end{enumerate}

\textbf{Construction.}
Let $M_0, M_1, M_2, \ldots$ be an effective enumeration of all
$\alpha$-bounded learners (i.e., all Turing machines equipped with
an $\alpha$-counter).  We define, for each $e \in \N$, a function
$f_e \in \mathcal{C}_\alpha$ designed to defeat learner $M_e$.

The function $f_e$ is defined by stages.  At each stage~$s$, $f_e$ has
been defined on arguments $0, 1, \ldots, s-1$.  We simulate $M_e$ on
the text $(f_e(0), f_e(1), \ldots, f_e(s-1))$ and observe $M_e$'s
hypothesis and counter value.  There are two cases:

\begin{itemize}[leftmargin=*]
\item \textbf{Case 1:} $M_e$ has exhausted its counter (the counter has
  reached~$0$).  Then $M_e$ can no longer change its mind.  It is now
  committed to some hypothesis $h$.  We define $f_e$ on the remaining
  arguments to disagree with $\varphi_h$: choose $f_e(s)$ so that
  $f_e(s) \neq \varphi_h(s)$ (this is possible since $\varphi_h$ is a
  fixed computable function and we are free to define $f_e$).  Then
  $M_e$ fails to $\Ex$-identify $f_e$.

\item \textbf{Case 2:} $M_e$ still has counter $> 0$.  We extend $f_e$
  by one more value, choosing $f_e(s)$ to force $M_e$ to change its mind
  (if possible) or simply setting $f_e(s) = 0$ if $M_e$ does not change
  its mind on any extension.
\end{itemize}

\textbf{Defeating $M_e$.}
By iterating Case~2, we either:
\begin{itemize}[leftmargin=*]
\item force $M_e$ to exhaust its $\alpha$-counter (each forced mind
  change decrements the counter, and there is no infinite descending
  chain in the ordinals), after which Case~1 applies; or
\item reach a point where $M_e$ refuses to change its mind no matter
  how we extend $f_e$, in which case $M_e$ has committed to a
  hypothesis $h$ and we can diagonalize as in Case~1.
\end{itemize}

In both sub-cases, $M_e$ fails on $f_e$.  Hence no $\alpha$-bounded
learner identifies all of $\mathcal{C}_\alpha = \{f_e : e \in \N\}$.

\medskip
\textbf{An $(\alpha+1)$-bounded learner for $\mathcal{C}_\alpha$.}
An $(\alpha+1)$-bounded learner $M^*$ for $\mathcal{C}_\alpha$ works
as follows.  Given text for some $f_e$, the learner $M^*$ searches
for the index $e$ by dovetailing: it simulates the construction of
$f_0, f_1, f_2, \ldots$ in parallel with the incoming data, and
hypothesizes the smallest $e$ consistent with the data seen so far.
Each time the current hypothesis $e$ is refuted by new data, $M^*$
moves to the next candidate, decrementing its counter.

The counter is initialized to $\alpha + 1$.  The first mind change
(from the initial hypothesis to the first genuine candidate)
decrements the counter from $\alpha + 1$ to $\alpha$.  That the
subsequent mind changes are bounded by~$\alpha$ is the crux of the
construction, and it is exactly here that this proof is a sketch: the
$\alpha$-bound on $M^*$'s residual revisions follows by reconstructing
the diagonal construction from the data via the recursion theorem, a
priority argument we do not carry out.  See Freivalds and
Smith~\cite{FreivaldsSmith1993} for the complete construction.
\end{proof}

\begin{proof}[Proof of \Cref{thm:ordinal-hierarchy}]
Inclusions hold by \Cref{lem:inclusion}.  For strictness: given
$\alpha < \beta$, the class $\mathcal{C}_\alpha$ from
\Cref{lem:witness} lies in $\Ex_{\alpha+1} \subseteq \Ex_\beta$ but
not in $\Ex_\alpha$.

For the final strict inclusion $\bigcup_\alpha \Ex_\alpha \subsetneq
\Ex$: this requires a class $\mathcal{C}$ that is $\Ex$-learnable but
not $\Ex_\alpha$-learnable for any constructive $\alpha$.  The
construction uses a priority argument: define $\mathcal{C}$ by
simultaneously diagonalizing against all $\alpha$-bounded learners
for all $\alpha$.  The learner for $\mathcal{C}$ is an ordinary
$\Ex$-learner with no ordinal counter; its mind-change count is finite
on each input but grows faster than any constructive ordinal as a
function of the input.  We refer to Ambainis, Jain, and
Sharma~\cite{AmbainisJainSharma1999} for the full priority construction,
which uses techniques from $\alpha$-recursion theory beyond the scope
of this chapter.
\end{proof}


\subsection{The hierarchy at limit ordinals}

The successor case $\Ex_\alpha \subsetneq \Ex_{\alpha+1}$ is the
engine of the hierarchy.  But limit ordinals also mark strict jumps.

\begin{prop}
\label{prop:limit-ordinals}
For every limit ordinal $\lambda$ with a notation in $\mathcal{O}$,
\[
  \bigcup_{\alpha < \lambda} \Ex_\alpha \subsetneq \Ex_\lambda.
\]
\end{prop}

\begin{proof}
The inclusion $\bigcup_{\alpha < \lambda} \Ex_\alpha \subseteq
\Ex_\lambda$ follows from \Cref{lem:inclusion}.  For strictness,
let $\lambda = \sup_n \alpha_n$ where $(\alpha_n)$ is the computable
sequence witnessing $\lambda \in \mathcal{O}$.

Define a class $\mathcal{C}_\lambda = \{f_n : n \in \N\}$ where
$f_n$ requires exactly $\alpha_n$ mind changes to identify.  (Each
$f_n$ is drawn from the witness class $\mathcal{C}_{\alpha_n}$ of
\Cref{lem:witness}.)  Then $\mathcal{C}_\lambda$ is not in
$\Ex_{\alpha_k}$ for any $k$ (since $f_{k+1}$ requires more than
$\alpha_k$ mind changes), but an $\Ex_\lambda$-learner can handle
it: upon seeing enough data to determine which $f_n$ is the target,
the learner sets its counter to $\alpha_n$ (which is
$<_{\mathcal{O}} \lambda$) and proceeds.
\end{proof}

\begin{remark}
\label{rem:omega-to-n}
The case $\lambda = \omega$ is particularly instructive.
$\Ex_\omega \setminus \bigcup_n \Ex_n$ consists of classes where the
learner needs finitely many mind changes on each input, but the
number of mind changes must be chosen adaptively.  The $\omega$-counter
captures exactly this deferred-commitment pattern.
\end{remark}


% ================================================================
% SECTION 4: ANOMALOUS LEARNING (brief)
% ================================================================

\section{Anomalous Learning}
\label{sec:anomalous}

The mind-change hierarchy refines $\Ex$ by bounding the number of
hypothesis revisions.  A different axis of relaxation tolerates
\emph{errors} in the final hypothesis.

\begin{defn}[$\Ex_a^*$ -- Anomalous Learning {\cite{CaseSmith1983}}]
\label{def:anomalous-ord}
For $a \in \N$, the class $\Ex_a^*$ consists of all concept classes
$\mathcal{C}$ for which there exists a learner that converges to a
hypothesis correct on all but at most $a$ inputs.  $\Ex^*$ allows
finitely many anomalies without a fixed bound:
$\Ex^* = \bigcup_{a \in \N} \Ex_a^*$.
\end{defn}

\begin{thm}[Anomaly Hierarchy {\cite{CaseSmith1983}}]
\label{thm:anomaly-hierarchy}
\[
  \Ex = \Ex_0^* \subsetneq \Ex_1^* \subsetneq \Ex_2^* \subsetneq \cdots
  \subsetneq \Ex^* \subsetneq \BC_0 \subsetneq \BC.
\]
Each inclusion is strict.
\end{thm}

\begin{remark}
\label{rem:anomaly-ordinal}
The anomaly hierarchy and the mind-change hierarchy are \emph{orthogonal}
axes of relaxation.  One can combine them: $\Ex_\alpha^a$ denotes
learning with an $\alpha$-bounded mind-change counter and tolerance for
$a$ anomalies.  The resulting two-dimensional hierarchy was studied by
Sharma, Stephan, and Ventsov~\cite{SharmaStephanVentsov2004}, who
showed that it yields a strict partial order under inclusion.
\end{remark}


% ================================================================
% SECTION 5: BC LEARNING (Profile D -- negative result)
% ================================================================

\section{Behaviorally Correct Learning and the Separation $\BC \setminus \Ex$}
\label{sec:bc-separation}

In $\Ex$-learning, convergence is \emph{syntactic}: the learner must
eventually output the same hypothesis index forever.  $\BC$-learning
relaxes this to \emph{semantic} convergence: the learner must eventually
output hypotheses that are all extensionally correct, but the index
may keep changing.

\begin{defn}[$\BC$-Learning {\cite{CaseSmith1983}}]
\label{def:bc-learning-ord}
A learner $M$ \emph{$\BC$-identifies} a function $f$ if there exists
$t_0$ such that for all $t \geq t_0$, $\varphi_{h_t} = f$
(extensional equality).  The hypothesis index $h_t$ may continue
to change.  The class $\BC$ consists of all concept classes
$\BC$-identifiable by some learner.
\end{defn}

The separation $\BC \setminus \Ex \neq \emptyset$ is a \emph{negative
result about $\Ex$}: there are classes where semantic convergence is
achievable but syntactic convergence is not.  The witness construction
is the content.

\begin{thm}[Case--Smith 1983]
\label{thm:bc-separation}
$\BC \setminus \Ex \neq \emptyset$.  There exists a class of total
recursive functions that is $\BC$-identifiable but not
$\Ex$-identifiable.
\end{thm}

\begin{proof}
\textbf{The witness class.}
For each total recursive function $g$ and each index $e$ with
$\varphi_e = g$, define the function $f_{g,e} : \N \to \N$ by
\[
  f_{g,e}(x) \;=\;
  \begin{cases}
    e & \text{if } x = 0, \\
    g(x) & \text{if } x \geq 1.
  \end{cases}
\]
The class is
\[
  \mathcal{C} \;=\; \bigl\{f_{g,e} : g \text{ total recursive},\;
  \varphi_e = g\bigr\}.
\]
Each total recursive function $g$ has infinitely many indices, so
$\mathcal{C}$ contains infinitely many functions that agree on
$\{1, 2, 3, \ldots\}$ but differ at argument~$0$.

\medskip
\textbf{$\mathcal{C} \in \BC$.}
A $\BC$-learner for $\mathcal{C}$ proceeds as follows.  Given text
for some $f_{g,e} \in \mathcal{C}$, the learner eventually sees
$f_{g,e}(0) = e$.  At time $t$, having observed
$f_{g,e}(0), f_{g,e}(1), \ldots, f_{g,e}(t)$, the learner outputs
an index $h_t$ for the function that agrees with the observed data
on $\{0, \ldots, t\}$ and follows $\varphi_e$ on $\{t+1, t+2, \ldots\}$.

Since $\varphi_e = g$ and $f_{g,e}(x) = g(x)$ for $x \geq 1$, each
hypothesis $h_t$ (for $t$ large enough that $f_{g,e}(0)$ has been
observed) computes $f_{g,e}$.  The indices $h_t$ may differ (the
learner may construct different programs at each step), but they all
compute the same function.  This is $\BC$-identification.

\medskip
\textbf{$\mathcal{C} \notin \Ex$.}
Suppose for contradiction that learner $M$ $\Ex$-identifies
$\mathcal{C}$.  Then for every $f_{g,e} \in \mathcal{C}$, $M$ must
converge to a single index $h^*$ with $\varphi_{h^*} = f_{g,e}$.

Fix any total recursive $g$.  For each index $e$ with $\varphi_e = g$,
the function $f_{g,e}$ belongs to $\mathcal{C}$.  The learner $M$,
on input $f_{g,e}$, sees $(e, g(1), g(2), g(3), \ldots)$ and must
converge to some $h_e^*$ with $\varphi_{h_e^*} = f_{g,e}$.

Consider two distinct indices $e \neq e'$ with $\varphi_e = \varphi_{e'} = g$.
The functions $f_{g,e}$ and $f_{g,e'}$ differ only at argument~$0$:
$f_{g,e}(0) = e \neq e' = f_{g,e'}(0)$.  So $M$ must converge to
distinct hypotheses $h_e^*$ and $h_{e'}^*$ (since
$\varphi_{h_e^*}(0) = e \neq e' = \varphi_{h_{e'}^*}(0)$).

But after time~$0$, both data streams are identical: $g(1), g(2), \ldots$.
The learner's behavior from time~$1$ onward depends only on its
hypothesis after seeing $f(0)$ and the shared tail.  Since $g$ has
infinitely many indices $e_0, e_1, e_2, \ldots$, the learner $M$ must
produce infinitely many distinct convergent hypotheses $h_{e_0}^*,
h_{e_1}^*, h_{e_2}^*, \ldots$, all determined by the single initial
datum $f(0)$ and the same infinite tail.

The contradiction is now information-theoretic, sharpened by the
recursion theorem.  Fix any total recursive~$g$ and let
$e_0, e_1, e_2, \ldots$ be infinitely many indices with
$\varphi_{e_i} = g$.  The inputs $f_{g,e_0}, f_{g,e_1}, \ldots$ share
the identical tail $g(1), g(2), \ldots$ and differ only in the initial
label.  After processing that label, $M$'s subsequent behavior is driven
by the \emph{same} tail, so it follows a single hypothesis trajectory;
yet $\Ex$-identification demands a \emph{distinct} final index
$h_{e_i}^*$ for each~$i$, since $\varphi_{h_{e_i}^*}(0) = e_i$.  A single
convergent trajectory cannot end at infinitely many distinct indices.
To turn this tension into a concrete failure, the operator recursion
theorem~\cite{CaseSmith1983} supplies a total recursive~$g$ together with
two of its indices $e \neq e'$ on which $M$ produces the \emph{same}
limiting hypothesis~$h^*$; then $\varphi_{h^*}(0)$ cannot equal both
$e$ and~$e'$, so $M$ misidentifies at least one of
$f_{g,e}, f_{g,e'}$.  We do not reproduce the priority bookkeeping of the
operator recursion theorem here.
\end{proof}

\begin{separation}
The class $\mathcal{C}$ witnesses $\BC \setminus \Ex \neq \emptyset$.
Here is why the $\Ex$-learner cannot escape.

A $\BC$-learner has room to maneuver: it may output different indices
for the same function, exploiting the many-to-one nature of G\"odel
numbering.  An $\Ex$-learner has no such room.  It must converge to a
\emph{single} index.

Now fill $\mathcal{C}$ with infinitely many functions that differ only
in their self-referential label at argument~$0$.  The $\BC$-learner
reads the label, builds a correct program, moves on.  The $\Ex$-learner
reads the label too, but it must commit.  Commit to which index?
Every candidate is correct extensionally for the current function but
wrong syntactically for all the others.  The learner converges; the
class shifts; the index is wrong.  The trap is shut.

The obstruction is the gap between extensional and intensional identity:
$\BC$ asks only that the program compute the right function; $\Ex$
demands the right program.
\end{separation}


% ================================================================
% SECTION 6: THE FULL LANDSCAPE (diagram)
% ================================================================

\section{The Full Landscape}
\label{sec:ordinal-landscape}

\Cref{fig:ordinal-hierarchy} displays the complete hierarchy of
mind-change-bounded and anomaly-bounded learning criteria.

\begin{figure}[htbp]
\centering
\begin{tikzpicture}[
  node distance=0.6cm and 1.2cm,
  every node/.style={font=\small},
  basecrit/.style={draw, very thick, rounded corners, minimum width=1.9cm,
    minimum height=0.58cm, align=center},
  litlevel/.style={draw, thin, dashed, rounded corners, minimum width=1.8cm,
    minimum height=0.54cm, align=center},
  litincl/.style={-{Stealth[length=5pt]}, thick, dashed, color=thmred},
  citetag/.style={font=\scriptsize\itshape, color=notegray},
]
% --- mind-change column (left, bottom to top): SOLID base FIN/Ex, DASHED ordinal levels ---
\node[basecrit, fill=defblue!10] (fin) {$\FIN = \Ex_0$};
\node[litlevel, fill=defblue!8, above=of fin] (ex1) {$\Ex_1$};
\node[litlevel, fill=defblue!8, above=of ex1] (ex2) {$\Ex_2$};
\node[above=0.26cm of ex2] (dots1) {$\vdots$};
\node[litlevel, fill=defblue!8, above=0.26cm of dots1] (exn) {$\Ex_n$};
\node[above=0.26cm of exn] (dots2) {$\vdots$};
\node[litlevel, fill=edgeorange!14, above=0.26cm of dots2] (exomega) {$\Ex_\omega$};
\node[litlevel, fill=edgeorange!14, above=of exomega] (exomega2) {$\Ex_{\omega \cdot 2}$};
\node[above=0.26cm of exomega2] (dots3) {$\vdots$};
\node[litlevel, fill=edgeorange!14, above=0.26cm of dots3] (exomega2a) {$\Ex_{\omega^2}$};
\node[above=0.26cm of exomega2a] (dots4) {$\vdots$};
\node[litlevel, fill=thmred!9, above=0.26cm of dots4] (exall) {$\bigcup_\alpha \Ex_\alpha$};
\node[basecrit, fill=thmred!14, above=of exall] (ex) {$\Ex$};
% --- anomaly arm (up-right from Ex): DASHED levels, SOLID BC ---
\node[litlevel, fill=sepgreen!10, right=2.4cm of ex] (exa1) {$\Ex_1^*$};
\node[litlevel, fill=sepgreen!10, above=of exa1] (exa2) {$\Ex_2^*$};
\node[above=0.26cm of exa2] (dots5) {$\vdots$};
\node[litlevel, fill=sepgreen!10, above=0.26cm of dots5] (exastar) {$\Ex^*$};
\node[litlevel, fill=analogypurple!10, above=of exastar] (bc0) {$\BC_0$};
\node[basecrit, fill=analogypurple!14, above=of bc0] (bc) {$\BC$};
% --- inclusion arrows: ALL dashed-literature (every rung is classical, none kernel-verified) ---
\draw[litincl] (fin) -- (ex1);
\draw[litincl] (ex1) -- (ex2);
\draw[litincl] (exn.north) -- (dots2.south);
\draw[litincl] (dots2.north) -- (exomega.south);
\draw[litincl] (exomega) -- (exomega2);
\draw[litincl] (exomega2a.north) -- (dots4.south);
\draw[litincl] (dots4.north) -- (exall.south);
\draw[litincl] (exall) -- (ex);
\draw[litincl] (ex) -- (exa1);
\draw[litincl] (exa1) -- (exa2);
\draw[litincl] (exastar) -- (bc0);
\draw[litincl] (bc0) -- (bc);
% --- one cite tag per author-result, at the regime boundary ---
\node[citetag, right=2pt of ex1] {Velauthapillai};
\node[citetag, left=2pt of exomega] {Freivalds--Smith};
\node[citetag, left=2pt of exall] {Ambainis--Jain--Sharma};
\node[citetag, right=2pt of exa2] {Case--Smith};
\node[font=\scriptsize, color=edgeorange, right=2pt of exomega2] {$\supsetneq \bigcup_n \Ex_n$};
% --- braces ---
\draw[decorate, decoration={brace, amplitude=6pt, mirror}, color=defblue!70]
  ([xshift=-1.7cm]fin.south west) -- ([xshift=-1.7cm]ex.north west)
  node[midway, left=8pt, align=center, font=\scriptsize, color=defblue!70] {mind-change\\hierarchy};
\draw[decorate, decoration={brace, amplitude=6pt}, color=sepgreen!70]
  ([xshift=1.6cm]exa1.south east) -- ([xshift=1.6cm]bc.north east)
  node[midway, right=8pt, align=center, font=\scriptsize, color=sepgreen!70] {anomaly\\hierarchy};
% --- VERIFIED ordinal-rank inset: the kernel's actual transfinite contribution, a SEPARATE measure ---
\node[draw, very thick, rounded corners, fill=defblue!5, align=left, font=\scriptsize,
  text width=6.1cm, below=1.4cm of fin.south west, anchor=north west] (rankbox) {%
  \textbf{Verified ordinal rank} (dimension side, a separate measure)\\
  \leanname{withTopNatToOrdinal} embeds $\N\cup\{\infty\}$ into the ordinals, sending $\infty$ to $\omega$\\
  \leanname{vcdim_to_ordinal_vcdim}: finite VC gives a finite rank, $\VC{=}\infty$ gives $\omega$\\
  finite-capacity classes (softmax attention, halfspaces) sit at the bottom, below $\omega$};
% --- evidence-grade legend ---
\coordinate (leg) at ([yshift=-0.42cm]rankbox.south west);
\draw[very thick] (leg) -- ++(0.55,0);
\node[font=\scriptsize, anchor=west] at ([xshift=0.62cm]leg) {kernel-verified ($sc{=}0$)};
\draw[thick, dashed, color=thmred] ([xshift=4.0cm]leg) -- ++(0.55,0);
\node[font=\scriptsize, anchor=west] at ([xshift=4.62cm]leg) {classical literature (cited)};
\draw[thick, dashed, color=edgeorange] ([yshift=-0.42cm]leg) -- ++(0.55,0);
\node[font=\scriptsize, anchor=west] at ([xshift=0.62cm, yshift=-0.42cm]leg)
  {open: not yet machine-checked (\cref{op:formalize-ordinal-hierarchy}, \cref{op:formalize-bc})};
\end{tikzpicture}
\caption[The ordinal complexity spine of limit identification: a classical tower with a verified rank at its base.]{The
number of mind changes a limit learner needs (left column) or the number of anomalies it
tolerates (right column) is a constructive ordinal, and that ordinal stratifies the
identification criteria into a strict tower: from finite identification $\FIN = \Ex_0$ up
through the transfinite levels $\Ex_\alpha$ into explanatory learning $\Ex$, and on through the
anomaly levels $\Ex_a^*$ into behaviorally correct learning $\BC$.  Every strict inclusion is a
classical theorem (the finite levels by Velauthapillai, the first limit jump
$\bigcup_n \Ex_n \subsetneq \Ex_\omega$ by Freivalds--Smith, the transfinite levels and
$\bigcup_\alpha \Ex_\alpha \subsetneq \Ex$ by Ambainis--Jain--Sharma, the anomaly axis by
Case--Smith), and each inclusion is drawn dashed because none is machine-checked and formalizing
the tower is still open.  The three base criteria $\FIN$, $\Ex$, $\BC$ are formalized, so their
nodes are drawn solid; the ordinal and anomaly levels are literature definitions and drawn
lighter.  The kernel certifies a separate ordinal rank on the dimension side, in the lower box:
an order-preserving embedding sends an infinite VC dimension to the first transfinite rank
$\omega$ and fixes every finite dimension below it, so every finite-capacity class of this book,
including softmax attention and the geometric classes, sits at a finite rank at the very bottom.
The transfinite tower thus becomes necessary only for the exact identification of programs.}
\label{fig:ordinal-hierarchy}
\end{figure}

We summarize the key structural facts.

\begin{thm}[Summary of Inclusions]
\label{thm:summary}
The following inclusions are all strict:
\begin{enumerate}[label=(\roman*)]
\item $\FIN = \Ex_0 \subsetneq \Ex_1 \subsetneq \cdots \subsetneq
  \Ex_n \subsetneq \cdots$ \quad (finite levels,
  Velauthapillai~\cite{Velauthapillai1989}).
\item $\bigcup_{n < \omega} \Ex_n \subsetneq \Ex_\omega$ \quad
  (first limit ordinal jump, Freivalds--Smith~\cite{FreivaldsSmith1993}).
\item $\forall \alpha < \beta$ constructive:
  $\Ex_\alpha \subsetneq \Ex_\beta$ \quad
  (\Cref{thm:ordinal-hierarchy}).
\item $\bigcup_{\alpha \in \mathcal{O}} \Ex_\alpha \subsetneq \Ex$
  \quad (Ambainis--Jain--Sharma~\cite{AmbainisJainSharma1999}).
\item $\Ex = \Ex_0^* \subsetneq \Ex_1^* \subsetneq \cdots \subsetneq
  \Ex^* \subsetneq \BC_0 \subsetneq \BC$ \quad
  (Case--Smith~\cite{CaseSmith1983}).
\end{enumerate}
\end{thm}


\subsection*{What the ordinal measures}

The mind-change ordinal of a class $\mathcal{C}$ is the least
constructive ordinal $\alpha$ such that $\mathcal{C} \in \Ex_\alpha$,
if such an ordinal exists.  If $\mathcal{C} \in \Ex \setminus
\bigcup_\alpha \Ex_\alpha$, then $\mathcal{C}$ has no ordinal
mind-change bound; its mind-change complexity exceeds all constructive
ordinals.

\begin{defn}[Mind-Change Ordinal of a Class]
\label{def:mco-class}
The \emph{mind-change ordinal} of a class $\mathcal{C}$ is
\[
  \mathrm{mco}(\mathcal{C}) \;=\;
  \min\bigl\{\alpha \in \mathcal{O} : \mathcal{C} \in \Ex_\alpha\bigr\}
\]
if the minimum exists, and $\infty$ otherwise.
\end{defn}

\begin{example}
\label{ex:pattern-languages}
Unions of $n$ pattern languages have mind-change complexity exactly
$\omega^n$ (Ambainis--Jain--Sharma~\cite{AmbainisJainSharma1999}, who
establish both the $\omega^n$ upper bound and a matching lower bound).
This provides a natural example of classes at each level $\omega^n$
in the hierarchy.
\end{example}

\begin{graphnode}[title={\nde{mind\_change\_ordinal}}]
\textbf{Graph node.} Category: \texttt{complexity\_measure}, Layer~5.\\
\textbf{Edges:}
$\nde{mind\_change\_ordinal}
  \xrightarrow{\rel{extends\_grammar}} \nde{mind\_change\_count}$
(new primitives required: constructive ordinals).\\
$\nde{mind\_change\_characterization}
  \xrightarrow{\rel{characterizes}} \nde{mind\_change\_ordinal}$
(Ambainis--Jain--Sharma 1999).\\
\textbf{Provenance:} Freivalds--Smith 1993.
\end{graphnode}

\begin{graphnode}[title={\nde{bc\_learning}}]
\textbf{Graph node.} Category: \texttt{success\_criterion}, Layer~4.\\
\textbf{Edges:}
$\nde{bc\_learning} \xrightarrow{\rel{restricts}} \nde{ex\_learning}$
(constraint removal: drops syntactic convergence requirement).\\
\textbf{Provenance:} Case--Smith 1983.
\end{graphnode}


\subsection{A Machine-Checked Ordinal Rank}
\label{sec:ordinal-rank-formal}

The mind-change ordinal is one ordinal-valued complexity measure, native
to recursion-theoretic identification.  The combinatorial dimensions of
\Cref{ch:dimensions} carry a second one, and that one has been formalized.

The Littlestone and VC dimensions take values in $\N \cup \{\infty\}$.
There is a machine-checked, order-preserving embedding of this value type
into the ordinals that fixes each finite dimension and sends the infinite
value to~$\omega$,\formal{FLT_Proofs.Bridge.withTopNatToOrdinal} and under
it the VC dimension is dominated by its ordinal lift,
$\mathrm{withTopNatToOrdinal}(\VC(\mathcal{C})) \leq
\mathrm{OrdinalVCDim}(\mathcal{C})$.\formal{FLT_Proofs.Bridge.vcdim_to_ordinal_vcdim}
This is the verified core of the ``ordinal VC dimension'' of Martin,
Sharma, and Stephan~\cite{MartinSharmaStephan2006}: a class of infinite VC
dimension sits at the first transfinite rank~$\omega$, finite-dimension
classes strictly below it.

\begin{prop}[Online mistakes descend an ordinal rank]
\label{prop:online-ordinal-descent}
In the online (Littlestone) model, the optimal mistake bound of a class
equals its Littlestone dimension.\formal{FLT_Proofs.Theorem.Online.optimal_mistake_bound_eq_ldim}
For a class of \emph{finite} Littlestone dimension, the Standard Optimal
Algorithm realizes this bound by treating the dimension as a rank: each
time it errs, the Littlestone dimension of the surviving version space
strictly decreases.  Lifted through the embedding above, this rank is a
finite ordinal, and the strict decrease on each mistake is precisely the
well-founded descent that \Cref{def:ordinal-mc} carries into the
transfinite.
\end{prop}

\begin{proof}[Proof idea]
The mistake-bound equality is an adversary argument matched against the
algorithm's potential.  Its engine is the machine-checked descent step:
a mistake on a version space of finite Littlestone dimension strictly
lowers that dimension, so iterating bounds the mistakes by the initial
rank.  The online counter never passes~$\omega$; it starts at a finite
rank (or, for an infinite-dimensional class, at~$\omega$) and counts down
by ordinary induction.  The transfinite levels $\Ex_{\omega^2},
\Ex_{\omega^\omega}, \ldots$ of this chapter are the price of
\emph{identifying programs} rather than \emph{predicting labels}; the
prediction game stays at the bottom of the ordinal tower.
\end{proof}

\begin{remark}[Attention sits at a finite rank]
\label{rem:attention-ordinal}
The softmax-attention concept class of \Cref{ch:generalization} is a
machine-checked well-behaved class,\formal{TLT.Bridge.softAttention_wellBehaved}
and \Cref{prop:certified-attention} gives it a finite Dudley capacity once
its parameters are norm-bounded.  Its ordinal rank is therefore finite:
attention lives at the very bottom of the ordinal picture, alongside
halfspaces and the other geometric classes.  The transfinite mind-change
hierarchy is invisible to such classes.  It becomes necessary only when
the objects learned are programs and the success criterion is exact
identification (the recursion-theoretic regime of this chapter), which no
finite-capacity concept class can reach.  Placing attention in the ordinal
landscape thus also marks where that landscape first becomes non-trivial.
\end{remark}


% ================================================================
% OPEN PROBLEMS
% ================================================================

\section{Open Problems}
\label{sec:ordinals-open}

The mind-change hierarchy is classically complete in its main lines, but
its formal status, its interaction with the verified online rank, and its
two-dimensional refinement remain open.  Each problem is tagged a
\emph{known unknown} (KU), sharply posed but unresolved, or an
\emph{unknown known} (UK), where the right object or invariant is still
missing; a design-lab or literature anchor is named where one exists.

\begin{openprob}[The mind-change/anomaly lattice]
\label{op:two-dim-lattice}
Does every point of the two-dimensional mind-change$\,\times\,$anomaly
lattice admit a natural characterization, as unions of $n$ pattern
languages characterize $\omega^n$?  A known unknown left open by Sharma,
Stephan, and Ventsov~\cite{SharmaStephanVentsov2004}.
\end{openprob}

\begin{openprob}[Formalizing the transfinite hierarchy]
\label{op:formalize-ordinal-hierarchy}
The dimension-to-ordinal bridge is machine-checked
(\leanname{withTopNatToOrdinal}), but the Freivalds--Smith hierarchy
itself (constructive ordinals, $\alpha$-bounded learners, the strict
inclusions of \Cref{thm:ordinal-hierarchy}) has no verified counterpart.
Can the ordinal mind-change hierarchy be formalized in a proof assistant,
including the priority constructions?  An unknown known: the recursion-
theoretic machinery has no formal home yet.
\end{openprob}

\begin{openprob}[Notation-invariant levels]
\label{op:notation-invariance}
\Cref{rem:notation-dependence} records that $\Ex_a$ can depend on the
notation~$a$, not just the ordinal~$|a|$.  Characterize exactly which
ordinals admit a notation-invariant $\Ex_\alpha$, and whether the
strictness theorem survives the quotient by notation.  A known
unknown~\cite{MartinSharmaStephan2006}.
\end{openprob}

\begin{openprob}[The ordinal rank of attention]
\label{op:attention-ordinal-rank}
\Cref{rem:attention-ordinal} places the attention class at a finite VC
rank, but its \emph{Littlestone} dimension (the online mind-change
rank) may still be infinite (thresholds already are).  Is the Littlestone
dimension of a softmax-attention class finite, and if so what is its exact
ordinal rank?  A known unknown on the transformer side
(\leanname{softAttention_wellBehaved}, \leanname{optimal_mistake_bound_eq_ldim}).
\end{openprob}

\begin{openprob}[A verified online-to-transfinite bridge]
\label{op:online-transfinite}
The online descent (each mistake strictly lowers the Littlestone
dimension) is machine-checked (\leanname{ldim_strict_decrease_on_mistake})
and stays below~$\omega$.  Is there a single formal framework in which this
verified finite descent and the transfinite mind-change counter are
instances of one well-founded-recursion principle?  An unknown known: the
unifying ordinal-recursion abstraction is not yet written down.
\end{openprob}

\begin{openprob}[Formalizing $\BC \setminus \Ex$]
\label{op:formalize-bc}
The separation of \Cref{thm:bc-separation} rests on the operator recursion
theorem, which we sketched rather than proved.  Can the Case--Smith
separation be machine-checked, operator recursion theorem
included~\cite{CaseSmith1983}?  A known unknown of formalization.
\end{openprob}

\begin{openprob}[Tightness of the Hessenberg bound]
\label{op:hessenberg-tight}
Exercise~3 gives $\mathcal{C}_1 \cup \mathcal{C}_2 \in
\Ex_{\alpha \oplus \beta}$ for the natural (Hessenberg) sum.  Is this bound
tight: are there classes whose union's mind-change ordinal is exactly
$\alpha \oplus \beta$ and no smaller?  A known unknown about resource
composition.
\end{openprob}

\begin{openprob}[Natural classes at each ordinal]
\label{op:natural-witnesses}
Pattern-language unions populate the levels $\omega^n$.  Is there a natural
family of concept classes populating the levels up to $\varepsilon_0$, or
does naturalness break down past $\omega^\omega$?  A known
unknown~\cite{AmbainisJainSharma1999}.
\end{openprob}

\begin{openprob}[An ordinal anomaly hierarchy]
\label{op:ordinal-anomaly}
The anomaly count~$a$ in $\Ex_a^*$ is a natural number.  Is there a
well-founded transfinite anomaly hierarchy (an $\Ex_\beta^*$ indexed by
ordinals), and does it interleave with the mind-change ordinals?  An
unknown known: the transfinite anomaly object has not been defined.
\end{openprob}

\begin{openprob}[Mind-change ordinals and ordinal VC dimension]
\label{op:mco-vs-ordinalvc}
The chapter's mind-change ordinal and the verified ordinal VC dimension
(\leanname{vcdim_to_ordinal_vcdim}) are two ordinal-valued measures on
overlapping domains.  For classes where both are defined, how are they
related: is one always a bound on the other?  A known
unknown~\cite{MartinSharmaStephan2006}.
\end{openprob}


% ================================================================
% EXERCISES
% ================================================================

\subsection*{Exercises}

\begin{enumerate}[leftmargin=*,itemsep=4pt]

\item \textbf{Finite mind-change hierarchy.}
  Prove directly (without ordinals) that $\Ex_n \subsetneq \Ex_{n+1}$
  for every $n \in \N$.  \emph{Hint:} Diagonalize against all learners
  with at most $n$ mind changes.

\item \textbf{$\omega$-counter mechanics.}
  Let $\mathcal{C} = \{f_n : n \in \N\}$ where $f_n$ is the
  characteristic function of $\{0, 1, \ldots, n\}$.  Show that
  $\mathcal{C} \in \Ex_\omega \setminus \bigcup_n \Ex_n$.

\item \textbf{Ordinal arithmetic of mind changes.}
  Show that if $\mathcal{C}_1 \in \Ex_\alpha$ and
  $\mathcal{C}_2 \in \Ex_\beta$, then
  $\mathcal{C}_1 \cup \mathcal{C}_2 \in \Ex_{\alpha \oplus \beta}$
  where $\oplus$ denotes natural (Hessenberg) sum.

\item \textbf{BC does not collapse.}
  Show that $\BC_0 \subsetneq \BC_1$: construct a class that is
  $\BC$-identifiable with at most one anomaly but not with zero
  anomalies.

\item \textbf{Anomaly vs.\ mind change.}
  Prove that $\Ex_1^* \not\subseteq \Ex_\omega$: tolerance for one
  anomaly cannot be simulated by any ordinal mind-change bound (without
  anomaly tolerance).  \emph{Hint:} Use a class where the anomalous
  input cannot be detected in finite time.

\item \textbf{Counter initialization matters.}
  Show that the learning power of $\Ex_\alpha$ depends on the
  \emph{notation} for $\alpha$, not just the ordinal.  That is,
  two different notations $a, a' \in \mathcal{O}$ with $|a| = |a'|$
  may yield different classes $\Ex_a \neq \Ex_{a'}$.
  \emph{Hint:} Different notations for $\omega$ give different
  computable sequences for decrementing.

\end{enumerate}

\subsection*{Notes and Further Reading}

The mind-change hierarchy was introduced by Freivalds and
Smith~\cite{FreivaldsSmith1993}, building on the finite mind-change
bounds of Velauthapillai~\cite{Velauthapillai1989} and the systematic
study of mind changes by Fulk, Jain, and
Osherson~\cite{FulkJainOsherson1995}.  The full characterization using
constructive ordinals, including the strictness of the hierarchy and
its non-exhaustion of $\Ex$, was completed by Ambainis, Jain, and
Sharma~\cite{AmbainisJainSharma1999}.

The four types of mind-change counters (constant, ordinal, linearly
ordered, partially ordered) were introduced by Sharma, Stephan, and
Ventsov~\cite{SharmaStephanVentsov2004}, who showed that these form a
strict hierarchy of learning power.  The connection between
mind-change ordinals and ordinal VC dimension was explored by Martin,
Sharma, and Stephan~\cite{MartinSharmaStephan2006}.

$\BC$-learning was introduced by Case and Smith~\cite{CaseSmith1983}.
The comprehensive treatment of learning criteria hierarchies is in
Jain, Osherson, Royer, and Sharma~\cite{JainOshersonRoyerSharma1999},
which remains the definitive reference for Gold-style learning theory.

The interaction between mind changes and anomalies produces a
two-dimensional lattice that is still not fully mapped.  Whether every
point in this lattice has a natural characterization (analogous to
pattern language unions for $\omega^n$) remains an open question.


\subsection*{Counterfactual exercises: generalizing Figure~\ref{fig:ordinal-hierarchy}}

Each exercise perturbs the tower.  Argue, in the figure's grammar (the number of mind changes or anomalies a limit learner needs is a constructive ordinal that stratifies the criteria into a strict tower from $\FIN = \Ex_0$ up through $\Ex$ and on into $\BC$; every strict inclusion is a classical theorem with no machine-checked counterpart, so every inclusion arrow is drawn dashed and cited to Velauthapillai, Freivalds--Smith, Ambainis--Jain--Sharma, or Case--Smith; the three base criteria $\FIN$, $\Ex$, $\BC$ are formalized and drawn solid, the ordinal and anomaly levels are literature definitions and drawn lighter; and the kernel's only verified transfinite content is a \emph{separate} ordinal rank on the dimension side, the lower box, that sends an infinite VC dimension to $\omega$ and places every finite-capacity class at a finite rank below it), that the perturbed picture subsumes the concrete one as an \emph{instance} (a criterion or rank placement the verified core or the literature exhibits), a \emph{boundary} (one inclusion or one node grade toggled so the evidence grade or the axis of relaxation changes), or a \emph{frontier} (a literature reading whose promotion to a verified rank is open).

\begin{enumerate}[leftmargin=*,itemsep=3pt]
\item \textbf{The base criteria are the figure's solid nodes.} Read the three boxes $\FIN$, $\Ex$, $\BC$ apart from the arrows between them. Show this is the figure's defining verified instance at the node level: these three are kernel definitions (\leanname{BCLearnable} and its $\FIN$ and $\Ex$ companions in the criterion library, the same base criteria formalized in \Cref{ch:gold}), so they are drawn with solid borders, whereas the levels $\Ex_\alpha$ (\cref{def:ex-alpha}) and $\Ex_a^*$ (\cref{def:anomalous-ord}) are literature definitions drawn lighter; the node grade is independent of the edge grade, and at the nodes alone the endpoints of the tower are certified while its rungs are not.

\item \textbf{The finite base levels are the integer instance of the ordinal spine.} Restrict the left column to $\FIN = \Ex_0 \subsetneq \Ex_1 \subsetneq \cdots \subsetneq \Ex_n$. Show this is the figure's most concrete instance of ``mind changes are an ordinal'': a learner bounded to $n$ revisions sits at the finite ordinal $n$, the strict inclusions hold (\cref{thm:summary}~(i), \cite{Velauthapillai1989}), and the arrows are dashed because even this integer fragment, the easiest in the tower, has no machine-checked counterpart; the whole transfinite picture is the closure of this fragment under limits.

\item \textbf{Attention sits at the bottom finite rank of the verified inset.} Read the boxed dimension-side rank, not the Gold tower. Show this is the instance the inset exists to display: the softmax-attention class is machine-checked well-behaved (\cref{rem:attention-ordinal}, \leanname{softAttention_wellBehaved}) and has finite Dudley capacity once norm-bounded (\cref{prop:certified-attention}), so its ordinal rank is finite and it lives at the very bottom of the picture beside halfspaces, while the embedding \leanname{withTopNatToOrdinal} reserves the first transfinite rank $\omega$ for an infinite VC dimension; the transfinite tower is invisible to every finite-capacity class.

\item \textbf{The online mistake count is the verified finite-rank shadow.} Read \cref{prop:online-ordinal-descent} as the inset's mechanism. Show this is the instance that certifies a descending rank: the optimal online mistake bound equals the Littlestone dimension (\leanname{optimal_mistake_bound_eq_ldim}), and each mistake strictly lowers the surviving version space's Littlestone dimension, a well-founded descent that stays below $\omega$; lifted through \leanname{withTopNatToOrdinal} this is a finite ordinal counted down by ordinary induction, the machine-checked finite shadow of the transfinite mind-change counter the tower draws above $\omega$.

\item \textbf{The verified rank and the mind-change ordinal are two measures, drawn apart on purpose.} Ask why the inset is fenced off from the left column rather than placed on it. Show this is the figure's principal boundary: the left column is the recursion-theoretic mind-change ordinal (\cref{def:mco-class}), a literature object, while the inset is the dimension-side ordinal VC rank (\leanname{vcdim_to_ordinal_vcdim}, \leanname{OrdinalVCDim}), the verified core of Martin--Sharma--Stephan \cite{MartinSharmaStephan2006}; they are two ordinal-valued measures on overlapping domains and whether one bounds the other is open (\cref{op:mco-vs-ordinalvc}), so collapsing them onto a single axis would assert a comparison the kernel does not certify.

\item \textbf{The first limit jump is the boundary between integers and ordinals.} Toggle the rung $\bigcup_{n<\omega}\Ex_n \subsetneq \Ex_\omega$. Show this is the boundary where the spine becomes genuinely transfinite: a learner may need finitely many mind changes on each input with no uniform integer bound, so the supremum is $\omega$ and no $\Ex_n$ suffices (\cref{rem:omega-to-n}, \cref{thm:summary}~(ii), \cite{FreivaldsSmith1993}); the arrow is dashed like every other, and it is exactly this jump that the integer hierarchy of Exercise~2 cannot reach, which is why the ordinal index is forced.

\item \textbf{The anomaly axis is an orthogonal relaxation, not a higher rung.} Turn from the mind-change column to the right column $\Ex \subsetneq \Ex_1^* \subsetneq \cdots \subsetneq \Ex^* \subsetneq \BC_0 \subsetneq \BC$. Show this is the boundary between two axes of relaxation: the mind-change column bounds the number of hypothesis revisions, while the anomaly column tolerates errors in the final hypothesis and then drops syntactic for semantic convergence (\cref{thm:anomaly-hierarchy}, \cref{thm:bc-separation}, \cite{CaseSmith1983}); both columns are dashed literature, and the two axes are independent, so the figure draws them as separate columns rather than one chain.

\item \textbf{The tower is classically complete but entirely unformalized.} Ask which arrow in the figure is solid. Show this is the dashed-versus-solid boundary the figure is built to make legible: none of the inclusion arrows is solid, because the Freivalds--Smith hierarchy, including its priority constructions, has no verified counterpart (\cref{op:formalize-ordinal-hierarchy}) and the Case--Smith separation rests on an operator recursion theorem that is sketched rather than machine-checked (\cref{op:formalize-bc}, \cite{CaseSmith1983}); the tower is a complete body of classical mathematics whose formalization is open, which is the precise content of drawing every rung dashed.

\item \textbf{Formalizing the tower would recolor the whole left column.} Ask which single advance turns the dashed rungs solid. Show this is the frontier the dashed arrows mark: a proof-assistant development of constructive ordinal notations, $\alpha$-bounded learners, and the strict inclusions of \cref{thm:ordinal-hierarchy} (\cref{op:formalize-ordinal-hierarchy}), together with the operator recursion theorem behind $\BC \setminus \Ex$ (\cref{op:formalize-bc}), would promote the literature tower to a verified one and make its arrows solid; until then the contrast in the figure is exactly the contrast between the solid dimension-side inset and the dashed recursion-theoretic spine.

\item \textbf{Is the mind-change ordinal the rank of a well-founded version-space descent functor whose finite shadow is the verified online descent?} Read the left column and the inset as two instances of one principle. Show this is the figure's deepest frontier: the verified online descent lowers a finite Littlestone rank on each mistake (the mechanism behind \leanname{optimal_mistake_bound_eq_ldim}, \cref{prop:online-ordinal-descent}) and stays below $\omega$, while the mind-change counter descends a constructive ordinal that may exceed every $\omega^n$; whether both are instances of a single well-founded-recursion principle, a version-space descent functor whose ordinal rank is the mind-change ordinal and whose finite-rank restriction is the certified online descent, is open (\cref{op:online-transfinite}), as is whether every point of the two-dimensional mind-change-by-anomaly lattice admits such a rank (\cref{op:two-dim-lattice}, \cite{SharmaStephanVentsov2004}); proving and formalizing it would place the dashed tower and the solid inset under one verified abstraction, which is why this reading is drawn dashed.
\end{enumerate}
